\chapter{Introdução}

O cenário atual da gestão de saúde no Brasil é marcado por uma dicotomia: de um lado, sistemas legados robustos, porém burocráticos e de difícil usabilidade; de outro, soluções modernas em nuvem que, embora ágeis, sequestram a soberania dos dados do profissional médico e impõem custos operacionais crescentes.

O projeto \textbf{IntraClinica} surge como uma terceira via. Fundamentado no binômio \textit{``Medicine for Humans. Intelligence for Business.''}, o sistema busca devolver ao médico a autoridade sobre o seu fluxo de trabalho e sobre a informação gerada em sua prática.

\section{Objetivos}

Este relatório tem como objetivo principal detalhar a arquitetura e os processos do IntraClinica para a Diretoria Médica e parceiros estratégicos. Os pilares deste desenvolvimento são:

\begin{itemize}
    \item \textbf{Soberania Radical:} Garantir que a posse física e lógica dos dados resida onde o médico exerce sua vontade.
    \item \textbf{Excelência em Experiência Médica (Experiência Clínica Soberana (Ecossistema IntraClinica)):} Desenvolver interfaces que não concorram pela atenção do médico com o paciente.
    \item \textbf{Resiliência Operacional:} Implementar infraestrutura capaz de operar independente de instabilidades em serviços externos.
\end{itemize}

\section{Justificativa}

A necessidade de um sistema como o IntraClinica foi identificada através do diagnóstico de ``hemorragia financeira'' em clínicas de médio porte, onde o descontrole sobre o almoxarifado e a fadiga gerada por softwares hostis resultam em prejuízos diretos e queda na qualidade do atendimento clínico.
